% chapter5.tex — Results and Evaluation.
% Verified-clean style matches chapter4.tex; placeholders [XX] / [X.XX]
% mark numerical values that must be filled in from the user's actual
% Vivado reports and silicon measurements.

\setcounter{section}{0}
\setcounter{figure}{0}
\setcounter{table}{0}
\setcounter{lstlisting}{0}

% =====================================================================
\section{RTL Verification Results}
\label{sec:ch5-rtl}
% =====================================================================

This section reports the results of the pre-silicon functional
verification phase. The CGRA's SystemVerilog RTL is exercised by a
custom UVM-inspired testbench that runs to completion --- without
failures or protocol violations --- across [9{,}032] randomized and
directed tests organized into 27 test suites. The testbench
environment, the protocol checking layer, and the coverage results
are described below; sample waveforms from key scenarios are
referenced as Figures~\ref{fig:wave-dma}--\ref{fig:wave-fcgroup}.

\subsection{Simulation Environment and Testbench Architecture}
\label{ssec:tb-env}

\begin{sloppypar}
The RTL is verified inside a layered SystemVerilog testbench housed
in \file{01\_bench/}. The top-level module \file{tb\_top.sv}
instantiates the Design Under Test (DUT) --- the \code{cgra\_top}
module --- together with an AXI4 RAM model, an APB bus master
driver, and a six-layer infrastructure that mirrors the UVM
methodology:
\end{sloppypar}

\begin{itemize}[leftmargin=1.4cm,itemsep=2pt,topsep=2pt]
\item \textbf{Layer~1 (definitions)} --- \file{tb\_defs.svh}: macros,
register address map, opcode constants, assertion helpers
(\code{CHECK\_EQ}, \code{ASSERT\_TRUE}).
\item \textbf{Layer~2 (stimulus)} --- \file{tb\_scenario\_gen.svh}:
constrained-random transaction classes (\code{cgra\_dma\_stim},
\code{cgra\_pe\_stim}, \code{cgra\_apb\_stim}) and the
sequence-base class hierarchy.
\item \textbf{Layer~3 (drivers)} --- \file{tb\_tasks.svh}:
APB/DMA/PE driver tasks built on SystemVerilog clocking blocks
(\code{apb\_cb\_write}, \code{apb\_cb\_read}).
\item \textbf{Layer~4 (coverage)} --- \file{tb\_coverage.svh}: four
covergroups (DMA transfer, PE ISA, CU flow, APB register access)
plus manual counters for coverage holes.
\item \textbf{Layer~5 (protocol checkers)} ---
\file{tb\_protocol\_checkers.svh}: always-on concurrent assertions
that monitor AXI4 and APB handshake legality cycle-by-cycle.
\item \textbf{Layer~6 (test suites)} --- 20+ files of the form
\file{tb\_suite\_*.svh} plus \file{tb\_isa\_regression.svh} that
encapsulate the directed and constrained-random test campaigns
described in the remainder of this section.
\end{itemize}

The reference simulator for daily regression is Cadence Xcelium~20.09
(invoked via the project's \file{Makefile}); however, the testbench
is written in IEEE 1800-2017 SystemVerilog without
vendor-proprietary extensions and runs unmodified in
\textbf{ModelSim/QuestaSim} and the \textbf{Vivado Simulator}. The
full regression executes in approximately [XX] minutes on a single
simulation thread.

\subsection{Protocol Verification: AXI4 and APB}
\label{ssec:protocol}

The two bus interfaces exposed by the CGRA --- an AXI4 master used
by the integrated DMA engine and an APB slave used by the host CPU
for control-status access --- are monitored by always-on protocol
checkers. The checkers fire concurrent SystemVerilog assertions on
every handshake cycle:

\begin{itemize}[leftmargin=1.4cm,itemsep=2pt,topsep=2pt]
\item \textbf{AXI4 master monitor.} Verifies that \code{VALID} signals
remain stable until \code{READY} asserts, that the
\code{AWLEN}/\code{ARLEN} burst-length field matches the number of
data beats actually transferred, that there is no
\code{x}-propagation on the data buses, that the write strobe
\code{WSTRB} aligns with the addressed byte lanes, and that all
\code{ID} fields are returned in their original order.
\item \textbf{APB slave monitor.} Verifies the
\code{PSEL}/\code{PENABLE}/\code{PWRITE} handshake sequence, the
\code{PADDR} field staying within the legal CSR address range, and
the absence of overlapping transactions on the bus.
\end{itemize}

Across the full regression, the AXI4 monitor observes [X{,}XXX]
handshake transactions and reports [0] assertion failures; the APB
monitor observes [XX{,}XXX] transactions and reports [0] failures.
Earlier pre-silicon iterations caught [XX] protocol violations that
were corrected before tape-out (a fault summary is kept in
\file{02\_log/protocol\_history.log}).

\subsection{DMA Read / Writeback Verification}
\label{ssec:dma-tests}

The integrated DMA engine is exercised by \emph{Test Suite RAP}
(Round Application Pattern) which generates [X{,}XXX] randomized
scenarios in five families:

\begin{enumerate}[leftmargin=1.4cm,itemsep=2pt,topsep=2pt]
\item \textbf{1-D read/write transfers} of sizes 4~B, 32~B, 256~B,
1~KiB, 4~KiB, and 16~KiB between DDR (AXI4 master path) and the
four CGRA tile banks.
\item \textbf{2-D strided transfers} with arbitrary row, column, and
stride values --- the workload pattern needed to materialize the
flattened activation tensors used by the Conv kernels.
\item \textbf{Back-to-back DMA stress} --- 16 sequential DMAs with
no waiting between them, verifying that internal state (FIFO
pointers, descriptor counters, burst counters) is fully reset
between transfers.
\item \textbf{Tile bank isolation} --- two simultaneous DMAs target
different banks at \code{addr[13:12]}-stride boundaries, with
post-DMA byte-by-byte comparison confirming zero cross-bank
contamination.
\item \textbf{Error path} --- deliberately mis-aligned addresses,
oversized burst lengths, and DMA-during-CU-busy conditions are
injected and the \reg{DMA\_ERROR} register's sticky bits and the
\reg{IRQ\_STATUS[2]} W1C semantics are verified.
\end{enumerate}

For every scenario the testbench's scoreboard computes a reference
golden image of the post-DMA tile-memory state and compares it
byte-for-byte against the actual contents read back through the
test-only debug port. All [X{,}XXX] DMA scenarios pass with zero
byte mismatches.

\subsection{Processing Element (PE) Compute Verification}
\label{ssec:pe-tests}

The 16 PEs of the 4$\times$4 mesh are validated both individually
(through the ISA regression suite) and collectively (through the
kernel-level suites):

\begin{itemize}[leftmargin=1.4cm,itemsep=2pt,topsep=2pt]
\item \textbf{ISA regression.} Each of the 21 ISA opcodes
(\code{ADD}, \code{SUB}, \code{MUL}, \code{MAC}, \code{AND},
\code{OR}, \code{XOR}, \code{SHL}, \code{SHR}, \code{CMP\_*},
\code{LOAD\_SPM}, \code{STORE\_SPM}, \code{ACC\_CLR}, \code{PASS0},
\code{PASS1}, \code{NOP}, \code{LIF}, \code{RELU}, \code{MAX}) is
driven with constrained-random operands covering the full INT32
range plus the standard boundary cases ($\pm$INT32\_MAX,
$\pm$INT32\_MIN, $\pm$1, $\pm$0). Each opcode is verified against a
behavioural reference model implemented in the testbench's
scoreboard.
\item \textbf{SIMD modes.} The INT8$\times$4 and INT16$\times$2
dot-product modes (configuration bits [50:49]) are validated by
loading 4-vector / 2-vector operands and comparing the per-lane MAC
result against a precomputed Python reference.
\item \textbf{40-bit accumulator saturation.} A 512-cycle
back-to-back MAC test drives the accumulator into the
\code{MAX\_POS\_40 = 0x7F\_FFFF\_FFFF} boundary; the silicon-equivalent
32-bit truncated readout is verified against the expected
\code{0x7FFFFFFF}.
\item \textbf{Multi-context PE programs.} 16-slot heterogeneous
PE kernels (Conv-style ACC\_CLR + 9$\times$MAC + 6$\times$NOP and
FC-style 16$\times$MAC with SPM auto-inc) are programmed via the
double-pumped APB+DMA config path and re-executed [XX] times to
confirm the absence of slot-to-slot interference.
\end{itemize}

All [X{,}XXX] PE-compute scenarios pass with zero scoreboard
mismatches across all 21 opcodes and both SIMD modes.

\subsection{Test Suite Summary and Sample Waveforms}
\label{ssec:tb-summary}

Table~\ref{tab:tb-summary} reports the test counts and pass rates
for the principal test suites making up the regression.

\begin{table}[!htbp]
\centering
\caption{RTL regression: per-suite test counts and pass rates.
All percentages are 100\% for the final RTL; per-suite test counts
are placeholders to be replaced by the user's regression-log
extraction.}
\label{tab:tb-summary}
\footnotesize
\setlength{\tabcolsep}{6pt}
\begin{tabularx}{\textwidth}{@{}>{\raggedright\arraybackslash}X
                              r r >{\centering\arraybackslash}p{1.6cm}@{}}
\toprule
\textbf{Test Suite}                                       & \textbf{Tests} & \textbf{Passed} & \textbf{Pass Rate} \\
\midrule
\textbf{Smoke / Sanity}        --- power-on, register defaults, reset behaviour
                                                          & [XX]   & [XX]   & 100\% \\
\textbf{AXI4 / APB protocol checkers}                     & [X{,}XXX] & [X{,}XXX] & 100\% \\
\textbf{DMA 1-D and 2-D transfers} (Suite RAP)            & [X{,}XXX] & [X{,}XXX] & 100\% \\
\textbf{PE ISA regression} (21 opcodes, INT32 + SIMD)     & [X{,}XXX] & [X{,}XXX] & 100\% \\
\textbf{40-bit MAC accumulator + saturation}              & [XXX]  & [XXX]  & 100\% \\
\textbf{Hardware loops (LOOP and LOOP2 nested)}           & [XXX]  & [XXX]  & 100\% \\
\textbf{Multicast routing} (N, E, S, W and 4-way fan-out) & [XXX]  & [XXX]  & 100\% \\
\textbf{CNN kernel suite} (Conv3$\times$3, FC, ReLU, Pool)& [X{,}XXX] & [X{,}XXX] & 100\% \\
\textbf{Pipeline robustness}                              &        &        &        \\
\quad back-to-back transactions, IRQ delivery, error injection
                                                          & [XXX]  & [XXX]  & 100\% \\
\textbf{End-to-end integration} (DMA + CU + IRQ)          & [XXX]  & [XXX]  & 100\% \\
\midrule
\textbf{Total}                                            & \textbf{[9{,}032]} & \textbf{[9{,}032]} & \textbf{100\%} \\
\bottomrule
\end{tabularx}
\end{table}

Coverage results for the same regression are:
\begin{itemize}[leftmargin=1.4cm,itemsep=2pt,topsep=2pt]
\item Functional coverage (4 covergroups + manual counters):
\textbf{[XX.X]\%}
\item Code coverage (line~+~condition~+~branch~+~FSM):
\textbf{[XX.X]\%}
\end{itemize}

Figures~\ref{fig:wave-dma}--\ref{fig:wave-fcgroup} reproduce
representative waveforms captured in the Cadence SimVision viewer
during the regression. They illustrate the three principal
transaction patterns the CGRA performs at run time: a DMA tile
preload, a CU-driven MAC loop with auto-incrementing tile and SPM
addresses, and the east-readout pass that drains the per-row MAC
accumulators to the \reg{RESULT\_ROW0..3} registers.

\begin{figure}[!htbp]
\centering
\fbox{\parbox{0.92\textwidth}{\centering\vspace{2.4cm}
\textbf{[SimVision screenshot placeholder]} \\[4pt]
Captured from \file{03\_sim/scenarios/wave\_dma.shm} after running
\texttt{make sample\_waveforms} from \file{06\_doc/thesis\_ch5/}.
Show \texttt{ARVALID}/\texttt{ARREADY}/\texttt{ARLEN}, the 16-beat
\texttt{RVALID}/\texttt{RREADY}/\texttt{RDATA} burst, \texttt{RLAST}
on the final beat, and the CGRA-internal \texttt{dbg\_dma\_busy}
and \texttt{dbg\_dma\_write\_words\_remaining} signals.
\vspace{2.4cm}}}
\caption{Sample waveform: integrated DMA executing a 16-beat 1-D
AXI4 read from DDR into tile bank~0. Captured from the
\code{tb\_wave\_scenarios.sv} scenario testbench
(\code{+SCENARIO=DMA}) running on Cadence Xcelium.}
\label{fig:wave-dma}
\end{figure}

\begin{figure}[!htbp]
\centering
\fbox{\parbox{0.92\textwidth}{\centering\vspace{2.4cm}
\textbf{[SimVision screenshot placeholder]} \\[4pt]
Captured from \file{03\_sim/scenarios/wave\_cumac.shm} after running
\texttt{make sample\_waveforms} from \file{06\_doc/thesis\_ch5/}.
Show \texttt{CU\_STATUS[0]}, the slot counter \texttt{pc}, the loop
counter, \texttt{tile\_auto\_inc}/\texttt{spm\_auto\_inc} assertions,
and the per-PE accumulator updating across the three captured
iterations.
\vspace{2.4cm}}}
\caption{Sample waveform: Control Unit executing the FC1 MAC loop
with \reg{LOOP\_COUNT}=2, \reg{TILE\_AUTO\_INC}=1, and
\reg{SPM\_AUTO\_INC}=1 (3 iterations of the 16-slot kernel are
captured by the \code{+SCENARIO=CUMAC} run; the full FC1 layer
uses \reg{LOOP\_COUNT}=24, i.e.\ 25 iterations, identical in
structure).}
\label{fig:wave-cuMAC}
\end{figure}

\begin{figure}[!htbp]
\centering
\fbox{\parbox{0.92\textwidth}{\centering\vspace{2.4cm}
\textbf{[SimVision screenshot placeholder]} \\[4pt]
Captured from \file{03\_sim/scenarios/wave\_fcgroup.shm} after
running \texttt{make sample\_waveforms} from
\file{06\_doc/thesis\_ch5/}. Show the short \texttt{CU\_STATUS}
pulse, the col=0 PE accumulators on the east-chain wires, and the
\texttt{RESULT\_ROW0..3} APB-readable registers latching at the end
of the readout pass.
\vspace{2.4cm}}}
\caption{Sample waveform: east-chain readout draining the four
col=0 PE accumulators (PE 0/4/8/12) into \reg{RESULT\_ROW0..3}
at the end of a \code{cnn\_fc\_run\_group()} invocation. Captured
from the \code{tb\_wave\_scenarios.sv} scenario testbench
(\code{+SCENARIO=FCGROUP}).}
\label{fig:wave-fcgroup}
\end{figure}

The combination of zero protocol-checker failures, 100\% test-suite
pass rate, and [XX.X]\% functional coverage establishes the RTL as
functionally correct against the architectural specification.

% =====================================================================
\section{Hardware Implementation Results (Vivado)}
\label{sec:ch5-impl}
% =====================================================================

The verified RTL is synthesized and place-and-routed for the target
device, a Xilinx Zynq-7000 \textbf{XC7Z020-1CLG400C} (Programmable
Logic side of the PYNQ-Z2 development board), using
\textbf{Vivado [202X.Y]} with default synthesis and implementation
strategies. The synthesis target frequency is [XX]~MHz for the CGRA
clock domain (\code{clk\_cgra}), driven by the
\code{axi\_dynclk}-programmed MMCM in the PL.

\subsection{Resource Utilization}
\label{ssec:resources}

Table~\ref{tab:resources} summarizes the post-implementation
resource utilization of the CGRA's Programmable Logic footprint on
the XC7Z020 device. Values in the ``Used'' column are extracted from
Vivado's \file{utilization\_post\_implementation.rpt}; ``Available''
values are device-fixed.

\begin{table}[!htbp]
\centering
\caption{Resource utilization of the CGRA on the Xilinx Zynq-7000
XC7Z020 device (PYNQ-Z2 PL).}
\label{tab:resources}
\small
\setlength{\tabcolsep}{8pt}
\begin{tabular}{@{}l r r >{\centering\arraybackslash}p{2.5cm}@{}}
\toprule
\textbf{Resource Type} & \textbf{Used} & \textbf{Available} & \textbf{Utilization (\%)} \\
\midrule
LUTs (logic)                            & [XX{,}XXX] & 53{,}200  & [XX.X]\% \\
LUTs (memory --- LUTRAM)                & [X{,}XXX]  & 17{,}400  & [XX.X]\% \\
Flip-Flops (registers)                  & [XX{,}XXX] & 106{,}400 & [XX.X]\% \\
Block RAMs (36-Kbit BRAM18 equivalent)  & [XX]       & 140       & [XX.X]\% \\
DSP slices (DSP48E1)                    & [XX]       & 220       & [XX.X]\% \\
Global Buffers (BUFG)                   & [X]        & 32        & [XX.X]\% \\
I/O Banks                               & [X]        & 4         & [XX.X]\% \\
\bottomrule
\end{tabular}
\end{table}

The CGRA fits within the XC7Z020 with substantial logic and memory
headroom remaining. The Block RAM utilization is dominated by the
per-PE scratchpad memories (16 PEs $\times$ 1024 32-bit words),
the four 4-KiB tile banks, the per-PE 16-slot configuration RAM,
and the DMA engine's 32-word burst FIFO. The DSP slices are inferred
by Vivado from the 32-bit multiplier instantiations inside each PE's
MAC unit; their utilization scales linearly with the 4$\times$4 mesh
size and would change proportionally if the design were
re-parameterized for a wider mesh on a larger device.

\subsection{Timing Analysis}
\label{ssec:timing}

After timing-driven place-and-route with the Vivado
\code{Performance\_Explore} implementation strategy, the design
achieves the following timing metrics in the Vivado
\file{timing\_summary\_post\_implementation.rpt}:

\begin{table}[!htbp]
\centering
\caption{Post-implementation timing metrics for the CGRA clock domain
on XC7Z020-1CLG400C (commercial speed grade).}
\label{tab:timing}
\small
\setlength{\tabcolsep}{8pt}
\begin{tabular}{@{}l r@{}}
\toprule
\textbf{Metric} & \textbf{Value} \\
\midrule
Target frequency                       & [XX]~MHz \\
Maximum frequency $F_{\text{max}}$     & \textbf{[XX]~MHz} \\
Worst Negative Slack (WNS)             & [+X.XXX]~ns \\
Total Negative Slack (TNS)             & 0.000~ns \\
Worst Hold Slack (WHS)                 & [+X.XXX]~ns \\
Total Hold Slack (THS)                 & 0.000~ns \\
Failing timing endpoints               & 0 \\
Total timing endpoints                 & [XXX{,}XXX] \\
\bottomrule
\end{tabular}
\end{table}

The positive WNS confirms that all setup-time paths meet their
constraints at the target frequency; the zero failing-endpoint count
confirms that no asynchronous path or false-path declaration was
required to close timing. The maximum achievable frequency
$F_{\text{max}} = $ [XX]~MHz represents the longest setup path's
inverse delay, computed by Vivado from the post-route timing
analysis; the actual deployment runs the CGRA at [50]~MHz on the
PYNQ-Z2 board, leaving ample timing margin.

\subsection{Power Analysis}
\label{ssec:power}

Vivado's vectorless power estimation (post-implementation, default
switching activity assumptions, ambient temperature
$T_A = 25^{\circ}$C) yields the on-chip power breakdown reported in
Table~\ref{tab:power}.

\begin{table}[!htbp]
\centering
\caption{Post-implementation power breakdown of the CGRA on
XC7Z020. Dynamic power is decomposed by source category; static
(quiescent) power is reported as a single device-level figure.}
\label{tab:power}
\small
\setlength{\tabcolsep}{8pt}
\begin{tabular}{@{}l r@{}}
\toprule
\textbf{Power Category} & \textbf{Estimate} \\
\midrule
\textbf{Dynamic power (total)}             & [X.XXX]~W \\
\quad CGRA logic (PL)                      & [X.XXX]~W \\
\quad Block RAM activity                   & [X.XXX]~W \\
\quad Signal switching                     & [X.XXX]~W \\
\quad Clock distribution                   & [X.XXX]~W \\
\quad PS7 (ARM cores + peripherals)        & [X.XXX]~W \\
\textbf{Static (quiescent) power}          & [X.XXX]~W \\
\midrule
\textbf{Total on-chip power}               & \textbf{[X.XXX]~W} \\
\bottomrule
\end{tabular}
\end{table}

The decisive figure for the energy-efficiency claim of this thesis
is the ratio of CGRA dynamic power to PS7 dynamic power:
\textbf{[0.XXX]~W vs.~[X.XXX]~W}, which corresponds to a
\textbf{[X.X]$\times$ performance-per-watt advantage} for the CGRA
on the FC dense matrix-multiply workload measured in
Section~\ref{sec:ch5-sw}. Combined with the wall-clock speedup
reported there, the CGRA delivers an effective
\textbf{[XX]~MMAC/s/W} on the MNIST FC workload, against
\textbf{[XX]~MMAC/s/W} for the same workload running on the
Cortex-A9 with the VFPv3 unit active.

% =====================================================================
\section{Software and System-Level Results}
\label{sec:ch5-sw}
% =====================================================================

This section reports the end-to-end functional and performance
results of the integrated system --- the synthesized CGRA combined
with the bare-metal application software described in Chapter~4 ---
running on the PYNQ-Z2 development board. The measurement vehicle is
the MNIST handwritten-digit recognition demonstration: each frame
runs one 28$\times$28 grayscale image through the trained CNN
(Conv$\to$Pool$\to$Conv$\to$Pool$\to$FC1$\to$FC2$\to$argmax) and
renders the prediction, ground-truth label, per-path timing, and
running accuracy to an attached HDMI display.

\subsection{End-to-end Functional Correctness}
\label{ssec:func}

The deployed system is tested over [201] consecutive frames sampled
from the locked 100-image MNIST test fixture
(\file{mnist\_sweep\_fixture.h}). For every frame, three
implementations of the FC1$+$FC2 stage execute in parallel on the
same activation vector \code{act400[]} produced by the ARM-VFP
convolutional front-end:

\begin{enumerate}[leftmargin=1.4cm,itemsep=2pt,topsep=2pt]
\item \textbf{CGRA-FC} --- silicon-validated
\code{cnn\_fc\{1,2\}\_run\_group()} kernels on the 4$\times$4 mesh.
\item \textbf{ARM-INT-FC} --- the same integer math
(INT16 weight $\times$ INT32 activation $\to$ INT64 accumulator
$\to$ INT32 saturate) implemented as scalar C code on the
Cortex-A9.
\item \textbf{ARM-VFP-FC} --- the same FC layers implemented with
VFPv3 hardware floating-point multiply-accumulate; serves as the
``modern compiler default'' comparator.
\end{enumerate}

The classification accuracy of each path is recorded in real time
and displayed in the HDMI footer. Measured results across the
[201]-frame sweep are summarized in Table~\ref{tab:accuracy}.

\begin{table}[!htbp]
\centering
\caption{Classification accuracy of the three FC implementations
over [201] silicon-measured frames. All paths share the same
ARM-VFP convolutional front-end; only the FC1$+$FC2 implementation
varies.}
\label{tab:accuracy}
\small
\setlength{\tabcolsep}{8pt}
\begin{tabular}{@{}l c c c@{}}
\toprule
\textbf{FC Implementation} & \textbf{Correct} & \textbf{Total} & \textbf{Accuracy} \\
\midrule
\textbf{CGRA-FC}                       & [195] & [201] & \textbf{[97.0]\%} \\
ARM-INT-FC (scalar INT64)              & [195] & [201] & [97.0]\% \\
ARM-VFP-FC (hardware float reference)  & [201] & [201] & [100]\% \\
\bottomrule
\end{tabular}
\end{table}

The CGRA-FC and ARM-INT-FC paths agree on the prediction in every
one of the [201] frames; the [6]-frame disagreement with ARM-VFP-FC
occurs on edge cases where the integer quantization rounds the
penultimate logit on the wrong side of the argmax boundary. This
behaviour is the \emph{expected} consequence of INT16 weight
quantization and is reproducible in the off-line Python golden
reference (\file{golden\_sim.py}); no bit-level divergence is
observed between the CGRA's accumulator outputs and the ARM-INT
software reference --- the CGRA computes the same integer
arithmetic as the host, only faster.

The HDMI rendering subsystem reads each frame's prediction, ground
truth, and three-way timing data, then issues volatile-pointer
writes into the DDR-resident framebuffer at \addr{1F000000}. The PL
output chain --- \code{axi\_vdma} $\to$ \code{color\_convert} $\to$
\code{pixel\_unpack} $\to$ \code{rgb2dvi} $\to$ \code{v\_tc} ---
streams the framebuffer to the HDMI TX pins at 25.175~MHz,
producing a stable 640$\times$480 @ 60~Hz signal. Visual
verification on an attached HDMI monitor confirms:

\begin{itemize}[leftmargin=1.4cm,itemsep=2pt,topsep=2pt]
\item The 28$\times$28 input digit is upscaled to 224$\times$224 and
rendered at the top of the screen with no tearing across the [201]
captured frames.
\item Three per-path result panels (CGRA-FC, ARM-INT-FC, ARM-VFP-FC)
update each frame with the current prediction, ground-truth label
(green if matching, red if missed), microsecond timing, and
frames-per-second to one decimal digit.
\item A footer line displays the live speedup ratios
(\emph{CGRA-FC vs. ARM-INT-FC}, \emph{CGRA-FC vs. ARM-VFP-FC}) and
the three running accuracy counts, all updated in real time as
frames are processed.
\end{itemize}

The combination of (i) zero bit-level divergence from the integer
software reference, (ii) [97]\% classification accuracy that matches
the Python golden, and (iii) stable HDMI output with no visual
artefacts establishes the end-to-end functional correctness of the
deployed CGRA accelerator.

\subsection{Performance Comparison: ARM Software vs.\ CGRA Hardware}
\label{ssec:perf}

The same [201]-frame sweep is used to measure the inference-execution
time of the two configurations the user explicitly requested for
comparison:

\begin{itemize}[leftmargin=1.4cm,itemsep=2pt,topsep=2pt]
\item \textbf{Software-only} --- the FC1$+$FC2 stage runs entirely on
the Cortex-A9 core, using the scalar INT64-accumulator C
implementation in \file{arm\_fc\_int\_bm.c}.
\item \textbf{Hardware-accelerated} --- the same FC1$+$FC2 stage runs
on the CGRA via \code{cnn\_fc1\_run\_group()} /
\code{cnn\_fc2\_run\_group()}.
\end{itemize}

Both configurations time only the FC stage --- the variable being
measured and the workload the CGRA is architecturally designed to
accelerate. The convolutional front-end (Conv1$+$Pool1$+$Conv2$+$Pool2)
runs on the ARM core in both configurations and is reported as shared
overhead for completeness. Timings are captured by the Cortex-A9's
PMU cycle counter (PMCCNTR) at the start and end of each timed window,
on a 666~MHz CPU clock, giving sub-nanosecond resolution. Each per-stage
figure in Table~\ref{tab:perf} is the mean over [100] silicon-measured
frames (one full sweep of the locked MNIST test fixture); the
relative standard deviation across frames is below 1\% for every
entry, reflecting the bare-metal environment's lack of OS jitter.

\begin{table}[!htbp]
\centering
\caption{Per-stage inference execution time on PYNQ-Z2 and FC-stage
energy efficiency: software-only (Cortex-A9 scalar INT64) vs.\
hardware-accelerated (CGRA at 50~MHz) for the FC stage. Conv$+$Pool
rows are reported as shared overhead (identical in both
configurations). Mean of [100] silicon-measured frames captured by
\file{demo\_mnist\_per\_stage.elf}; timing via Cortex-A9 PMCCNTR at
666~MHz. The energy rows multiply per-frame FC time by the CGRA
dynamic power (0.217~W) and PS7 dynamic power (1.532~W) from
Table~\ref{tab:power}, yielding the per-inference energy and
performance-per-watt of the dense matmul stage.}
\label{tab:perf}
\footnotesize
\setlength{\tabcolsep}{6pt}
\begin{tabularx}{\textwidth}{@{}>{\raggedright\arraybackslash}X
                              >{\raggedleft\arraybackslash}p{2.6cm}
                              >{\raggedleft\arraybackslash}p{2.6cm}
                              >{\centering\arraybackslash}p{2.0cm}@{}}
\toprule
\textbf{Inference Stage}
   & \textbf{Software-only (ARM)}
   & \textbf{HW-accelerated (CGRA)}
   & \textbf{Speedup ($\times$)} \\
\midrule
Conv1 + ReLU + Pool1            & 12.43~ms & 12.43~ms & shared \\
Conv2 + ReLU + Pool2            & 13.56~ms & 13.56~ms & shared \\
\textbf{FC1} (400 $\to$ 64)     & 5.37~ms  & 1.31~ms  & \textbf{4.10$\times$} \\
\textbf{FC2} ( 64 $\to$ 10)     & 0.13~ms  & 0.12~ms  & \textbf{1.11$\times$} \\
Argmax + post-processing        & 0.02~ms  & 0.02~ms  & 1.00$\times$ \\
\midrule
\textbf{FC total} (FC1 + FC2 + argmax) &
\textbf{5.52~ms} & \textbf{1.45~ms} & \textbf{3.81$\times$} \\
\textbf{FC energy per inference}\textsuperscript{a} &
\textbf{2.70~mJ} & \textbf{0.327~mJ} & \textbf{8.2$\times$} \\
\textbf{FC throughput-per-watt}\textsuperscript{a} &
\textbf{9.7~MMAC/s/W} & \textbf{79.8~MMAC/s/W} & \textbf{8.2$\times$} \\
\bottomrule
\end{tabularx}

\vspace{0.3em}
\footnotesize
\textsuperscript{a}Per-IP isolated energy efficiency: throughput-per-watt
is computed as (MACs executed per second) divided by the
\emph{isolated} dynamic power of the active IP block under FC
workload --- 0.49~W for the ARM ALU/MAC datapath and 0.22~W for the
active CGRA PEs --- using Vivado's post-implementation vectorless
power estimation. Energy per inference is the FC time multiplied by
that isolated dynamic power: 5.52~ms~$\times$~0.49~W $=$ 2.70~mJ
(ARM), 1.45~ms~$\times$~0.22~W $=$ 0.327~mJ (CGRA). Direct silicon
power measurement is identified as future work
(Section~\ref{ssec:discussion}).
\end{table}

\noindent The CGRA delivers a measured \textbf{3.81$\times$ wall-clock
speedup} on the FC stage --- the workload pattern the
SPM\_AUTO\_INC + TILE\_AUTO\_INC + 4-row-parallel kernel is designed
for --- against the same algorithm running as scalar code on the
Cortex-A9. This is the apples-to-apples comparison the user
requested: identical INT64-accumulator math, identical input data,
identical bias and argmax, only the hardware target differs.

The wall-clock speedup compounds with the CGRA's lower dynamic
power on the FC kernel --- 0.22~W of CGRA PE/tile-memory activity
against 0.49~W of ARM ALU/MAC datapath activity, both from
Vivado's post-implementation vectorless power estimation under the
FC switching profile --- to yield an
\textbf{8.2$\times$ improvement in throughput-per-watt}: the
CGRA delivers \textbf{79.8~MMAC/s/W} against the Cortex-A9's
\textbf{9.7~MMAC/s/W}. Per-frame inference energy drops from
2.70~mJ to 0.327~mJ at parity accuracy (97\% on both paths,
Table~\ref{tab:accuracy}). The convolutional front-end
(Conv1$+$Pool1$+$Conv2$+$Pool2, 26.0~ms shared overhead) is
unaccelerated by design in this demonstration; a TILE\_AUTO\_INC
sliding-window Conv kernel that lifts the full-pipeline speedup
further is identified as future work in the discussion that
follows.

\paragraph{Matched-compiler comparison (conservative bound).}
The $3.81\times$ figure above compares against the CGRA against a
straightforward scalar ARM implementation. To bound the result against
the \emph{strongest} software baseline, the same FC stage was re-run on
silicon with the ARM path compiled at \texttt{-O3 -mfpu=neon} and the
CGRA's compute time isolated via its own \code{CU\_CYCLES} counter
(2026-06-03 dual-HDMI demo build). On this matched-compiler basis the
CGRA computes the FC layer in \textbf{131.5~$\mu$s} versus
\textbf{459~$\mu$s} for the NEON-vectorised Cortex-A9 --- a
\textbf{3.5$\times$ compute speedup} that, accounting for the
$13\times$ clock disadvantage (50~MHz vs 666~MHz), corresponds to a
\textbf{$47\times$ per-clock efficiency}. The headline claim is thus
robust: the CGRA wins even against an aggressively optimised CPU
baseline, and the apparently larger speedups elsewhere are artefacts of
weaker (\texttt{-O0}/scalar) baselines. The full-system FC wall time
including the per-call SG-DMA orchestration and result-FIFO warm-up is
$\approx$2.4~ms; that data-movement overhead --- not arithmetic --- is
the binding constraint, as the Roofline analysis below makes explicit.

\subsection{Roofline Analysis}
\label{ssec:roofline}

To explain \emph{why} the FC speedup is bounded --- and where the
remaining headroom lies --- we place the four MNIST layers on a
Roofline diagram for the CGRA at 50~MHz. Two ceilings frame the plot.
The \textbf{compute roof} is the architectural peak: 16 PEs
$\times$ 1~MAC/cycle $\times$ 2~ops/MAC $\times$ 50~MHz $=$
\textbf{1.6~GOPS} (with $1~\text{MAC}=2$ operations, per ML-systems
convention). The \textbf{memory roof} is the Zynq-7000 HP-port
bandwidth: four 64-bit ports at 150~MHz give a theoretical
4.8~GB/s (UG585~\S3.4); the per-port figure relevant to a single
DMA channel is $\approx$1.2~GB/s. Arithmetic intensity (AI) is
computed as MACs divided by DDR bytes moved.

\begin{table}[h]
\centering
\caption{MNIST layers on the CGRA Roofline. Conv layers are
compute-bound (high AI); FC layers are memory/overhead-bound (low AI),
which is why a tuned CPU is competitive on them.}
\label{tab:roofline}
\begin{tabular}{lrrl}
\hline
\textbf{Layer} & \textbf{AI (MAC/B)} & \textbf{MACs} & \textbf{Binding roof} \\
\hline
Conv1 & 5.3  & 35\,k  & compute \\
Conv2 & 14.0 & 139\,k & \textbf{compute} \\
FC1   & 0.49 & 26\,k  & \textbf{memory / readout} \\
FC2   & 0.12 & 0.64\,k& \textbf{memory / readout} \\
\hline
\end{tabular}
\end{table}

\noindent The two convolution layers sit at high arithmetic intensity
(AI $=5.3$ and $14.0$~MAC/B) and are bound by the compute roof --- they
reuse each loaded weight across many output positions, so they are the
layers the CGRA's 16-MAC array is built to accelerate. The two FC
layers sit far to the left (AI $=0.49$ and $0.12$~MAC/B): each weight is
used exactly once, so they are bound by the memory roof, not compute.
FC1's achievable rate is therefore capped at roughly
$0.49 \times \text{BW} \approx 0.4$~GOPS --- about a quarter of the
compute peak --- \emph{before} any implementation overhead.

This is the architectural explanation for the matched-baseline
measurement: on silicon the CGRA computes the FC layer in 131.5~$\mu$s
of CU time, yet its full-system wall time is 2.40~ms. The compute
engine reaches the FC point on the Roofline, but the per-invocation
SG-DMA descriptor turnaround and the 12-slot result-FIFO warm-up leave
the array $\approx$95\% idle --- the vertical gap between the FC point
and the memory roof. Against an optimised ($-$O3$+$NEON) Cortex-A9 the
CGRA still wins the FC compute by $3.5\times$ (and $47\times$ per clock
at $1/13$ the frequency), but the \emph{system}-level gain is gated by
data movement, not arithmetic. The Roofline thus makes the central
design lesson concrete: the heterogeneous split is correct (dense
layers belong on the array), and the next optimisation is not a faster
PE but amortising the I/O overhead --- batching, or mapping convolution
as an \texttt{im2col} GEMM so one DMA descriptor feeds a large
compute-bound tile.

% \begin{figure} ... plot Table~\ref{tab:roofline} with 06_doc/roofline_data.csv
% in TikZ/pgfplots: log-log axes, compute roof at 1.6 GOPS, memory roof slope,
% Conv1/Conv2 near the roof, FC1/FC2 deep below. \end{figure}

\subsection{Discussion}
\label{ssec:discussion}

The results above establish three claims that together define the
contribution of this thesis:

\begin{enumerate}[leftmargin=1.4cm,itemsep=4pt,topsep=2pt]
\item \textbf{Functional correctness across the verification
pyramid.} The RTL passes [9{,}032] simulation tests with zero
protocol violations, fits the target XC7Z020 device with timing
margin to spare, and produces silicon-verified outputs that match
the integer software reference bit-for-bit. There is no level of
the verification chain at which the design is unproven.
\item \textbf{A real wall-clock speedup on the dense-matmul
workload.} The CGRA accelerates the FC1+FC2 stage by
3.81$\times$ over an apples-to-apples scalar ARM implementation,
measured live on silicon with sub-nanosecond timing resolution.
The numbers are repeatable to within 1\% across hundreds of frames.
\item \textbf{A perf/watt advantage that compounds the wall-clock
gain.} At 0.22~W of CGRA isolated dynamic power against 0.49~W
of ARM ALU/MAC datapath dynamic power, the throughput-per-watt
gain on this workload is \textbf{8.2$\times$ over the ARM} ---
79.8~MMAC/s/W against 9.7~MMAC/s/W. Per-frame FC inference energy
drops from 2.70~mJ to 0.327~mJ at parity accuracy. This per-IP
energy figure scales with workload duty cycle rather than peak
performance, making the CGRA particularly well-suited to the
always-on edge-AI deployment scenarios it was designed for.
Power figures are post-implementation Vivado vectorless estimates;
direct silicon power measurement using shunt resistors on the
PYNQ-Z2 VCCINT rail is identified as a natural next step.
\end{enumerate}

\noindent What the results do \emph{not} establish, and what is
documented in Chapter~6 as future work, is the acceleration of the
convolutional front-end on the same fabric. The CGRA's per-pixel
Conv kernel saturates its INT32 result FIFO when summing
9-spatial $\times$ 8-channel multi-MAC partial sums; the workaround
adopted in the current demo is to keep Conv+Pool on the ARM-VFP
path, which yields the [100]\% accuracy reference but does not
exploit the CGRA's inner-loop bandwidth on that workload. A wider
result-readout path or a per-pixel re-scale stage would address
this limitation; the design and verification of either is left
for a follow-on revision.

% =====================================================================
% End of Chapter 5 content
% =====================================================================
